\begin{figure}[p]
\centering
\includegraphics[width=\textwidth]{supplementary/curated/original_selector.pdf}
\caption{\textbf{The original scalar initialization score predicts a first step but fails at prospective endpoint selection.} \textbf{A}, Four balanced leaf assignments (balanced 1--4; the number under each leaf is the input index at that leaf slot) and one comb, each at route budgets $K=1,2,4,8$: twenty candidates with prescribed cost. Filled nodes mark the $K$ feedback channels, one subtree each. \textbf{B}, Development, independent calibration, sealed choices and confirmation training. \textbf{C}, Final test half-mean-squared-error (MSE) plus cost, expressed as excess over the best trained candidate, for the feedback-only (green) and joint passive-transfer (blue) arms; the dashed line marks zero regret. Small points are the twenty whole-seed means; open diamonds and whiskers are means and 95\% seed-bootstrap intervals (10,000 draws), which are narrower than the diamond in the rank-only, fixed / max-budget and random-expectation rows. The fixed / max-budget baseline (the development choice, which is also the maximum budget) and the privileged rank-only baseline outperform the moment selector; random expectation averages all candidates uniformly. \textbf{D}, Route budget of the candidate chosen by the moment selector versus imposed rank; small points are the twenty seed-block means, open circles their mean, and the gray dashed line the retrospectively best candidate's budget, which equals the rank in every task of both arms. \textbf{E}, Mean within-task Spearman correlation of initialization utility with the actual first-step test-loss decrease and with the final test-loss ranking; small points are the twenty seed-block means, open circles and whiskers their mean and 95\% interval resampling whole seeds ($n=20$ seed blocks). Colours in \textbf{D,E} as in \textbf{C}. All 12,800 confirmation candidate fits and 38,400 checkpoints remain included. Every model has a freely trained dense weight matrix; passive transfer changes conditioning and cost but does not necessarily restrict representation. This experiment is distinct from the locally connected multi-affine construction and the separately frozen finite-horizon follow-up.}
\label{fig:si_original_selector}
\label{fig:supp_original_prospective_selection}
\end{figure}

\begin{figure}[p]
\centering
\includegraphics[width=\textwidth]{supplementary/curated/finite_horizon.pdf}
\caption{\textbf{Finite-horizon prediction improves selection but strong simple baselines remain.} \textbf{A,B}, Final costed-test-loss regret by imposed rank in the feedback-only and joint-transfer fixed-cache learners; regret is excess above the best trained candidate. \textbf{C}, Differences among strong selectors at each rank in the joint arm. The independently calibrated Gaussian SGD forecast improves over the original scalar utility; the full-batch and observed-context-count baselines remain competitive. All finite-horizon policies choose $K=r$ in this family, so improvement does not establish fine morphology selection. \textbf{D}, Median measured elapsed time on one central processing unit per twenty-candidate decision. Plug-in timings include calibration-model preparation, count-only time excludes shared calibration generation, and the pilot includes sixteen candidate updates; times measure these implementations, not hardware-independent complexity. \textbf{E}, Original scalar forecast versus final test half-mean-squared-error in the feedback-only arm. \textbf{F}, The calibration-derived Gaussian SGD forecast on the same cases and arm. Hexagons encode counts on a logarithmic color scale and dashed lines mark equality; the joint arm reproduces \textbf{E,F} and is retained in Source Data. The population-oracle forecast bias --- predicted minus exact population loss of the realized trained weights under fresh examples or finite-cache reuse --- is retained in Source Data: the fresh-example recurrence is exact under the stated Gaussian and conditional-linear assumptions, whereas cache reuse violates independence. Curves and whiskers in \textbf{A--D} are means and pointwise 95\% intervals from 10,000 whole-seed bootstrap draws over twenty seed blocks. All 25,600 candidate fits and both regimes are retained. These are linear-learning forecast tests, not a general nonlinear dendritic morphology law.}
\label{fig:si_finite_horizon}
\label{fig:supp_finite_horizon_selection}
\label{fig:supp_finite_horizon_prediction}
\end{figure}

\begin{figure}[p]
\centering
\includegraphics[width=\textwidth]{supplementary/curated/morphology_estimation.pdf}
\caption{\textbf{Finite calibration, adaptive construction and direct gradient learning.} \textbf{A}, Primary calibration: 256 labels, noise standard deviation (SD) 0.5; 192 observations for estimation/pilot fitting and 64 for pilot gating. The 255-term Walsh estimator receives no family labels, true support or coefficients. The pilot fits all twelve candidates for two sweeps. Choices are sealed; training restarts independently. \textbf{B}, Prespecified pooled fixed-baseline-minus-estimated and pilot-minus-estimated normalized mean squared error (NMSE). Whiskers are Bonferroni-adjusted 97.5\% paired bootstrap intervals; the dotted line marks the 0.01-NMSE meaningful-improvement margin. \textbf{C}, Excess clean-test NMSE above the best trained candidate across calibration sizes; solid curves use noise SD 0.5 and dashed curves no noise. \textbf{D}, Family-specific primary regret for estimated interactions, the development-fixed/estimated-rank baseline and the fitting pilot; pooled superiority is not uniform across families. \textbf{E}, Secondary adaptive trees trained with alternating least squares (ALS), versus the best trained twelve-candidate menu and a full-target construction oracle. Random-interaction errors remain. \textbf{F}, Elapsed time (median and interquartile range) on one CPU for eighty tasks: estimation/scoring of twelve trees versus pilot fitting/gating, excluding final training. \textbf{G,H}, Independent twenty-seed Adam cohort. \textbf{G}, Three structures and two credit rules in four non-isospectral families: filled symbols, exact credit (Adam 0.01); open, root broadcast (0.003). Label-noise floor: 0.0225; \textbf{E} uses clean targets. \textbf{H}, Fixed-minus-estimated error under exact credit and broadcast-minus-exact error on estimated trees use Bonferroni-adjusted 97.5\% intervals; estimated-minus-target-informed error under exact credit uses a pointwise 95\% interval. Target-informed construction is not a performance ceiling. Other whiskers: pointwise 95\% intervals, 10,000 whole-seed bootstrap draws ($n=20$); pooled statistics average families within seed. Both cohorts are independent of Supplementary Fig.~\ref{fig:si_oracle_profile_credit}.}
\label{fig:si_morphology_estimation}
\label{fig:supp_morphology_calibration}
\end{figure}
